\scanpage{303}
\mspage{219}

a) $\underline{\operatorname{Hom}}(J,C)_{/\varphi}$ est grde devant $\pi_0$.

b) Pour tout $j\in J$, le foncteur

\[
\underline{\operatorname{Hom}}(J,C)_{/\varphi}\longrightarrow C_{/\varphi(j)}\qquad\left|\begin{aligned}&\text{foncteur entre}\\&\text{catégories filtrantes}\\&\text{grdes devant }\pi_0\\&\text{en vertu de a)}\end{aligned}\right.\tag*{(4.9.2.1)}
\]

induit par le foncteur

\[
\left|\begin{aligned}\underline{\operatorname{Hom}}(J,C)&\longrightarrow C\\ F&\longmapsto F(j)\end{aligned}\right.,\tag*{(4.9.2.2)}
\]

est \uline{cofinal}.

Le résultat principal \deleted{du}\inserted{de} \unsure{la} présente section est le suivant :

\uline{Théorème} 4.9.3. Sous les conditions générales de 4.9.1 pour $C,\pi_0,J$ (notamment $\operatorname{card}\operatorname{réd}J\leq\pi_0$), le foncteur pl fidèle (4.9.3) est une équivalence de catégories dans chacun des deux cas suivants :

a) $C$ est stable par $\varinjlim$ filtrantes dénombrables.

\margin{Cas \unsure{commun} : il existe un ordinal limite $\alpha$, tel que 1) \deleted{Pour \unsure{tt entier} $\beta<\alpha$, \illegible{1}} $\operatorname{card}\deleted{\unsure{\beta}}\alpha\leq\unsure{\pi|\pi_0}$. 2) les limites inductives de type $I_\alpha=\{\beta\mid\beta<\alpha\}$ \unsure{st|sont} représentables dans $C$}% slanted

b) Il existe un cardinal $\pi'_0<\pi_0$, et un cardinal $\pi''_0>\pi'_0$ (NB OPS que $\pi''_0$ est le successeur de $\pi'_0$)\inserted{donc $\pi_{\!0}^{\unsure{'|''}}\leq\pi$}, tels que $C$ soit stable par $\deleted{\varinjlim}$limites inductives \unsure{grdes} devant $\pi'_0$, et de cardinal $\leq\pi''_0$.

\uline{Remarques} 4.9.4. 1) Ce théorème, cas a), contient le th. 4.4.13, partie b) (p.~95-96), où on suppose que $C$ est stable par limites inductives de cardinal $\leq\pi_0$, l'hypothèse a) de \unsure{4.9.3} étant considérablement plus faible.
% archive sheet number: 302
